\documentclass[12pt,a4paper]{article}
\usepackage[T1]{fontenc}
\usepackage[utf8]{inputenc}
\usepackage[brazil]{babel}
\usepackage{lmodern}
\usepackage[a4paper,margin=2.3cm]{geometry}
\usepackage{booktabs,pgfplots,hyperref}
\pgfplotsset{compat=1.18}
\hypersetup{colorlinks=true,linkcolor=blue!50!black}
\setlength{\parskip}{0.4em}
\setlength{\parindent}{0pt}
\title{\textbf{Análise dos resultados do experimento ATLAS}\\[0.4em]\large Detecção de problemas arquiteturais}
\author{}
\date{Outubro de 2026}
\begin{document}
\maketitle

\begin{abstract}
Este relatório compara três modelos na identificação de problemas arquiteturais em dez códigos. Foram analisadas cinco repetições por código e modelo, totalizando 150 execuções. O GPT-5.6 Luna obteve F1 médio de 0,860, seguido pelo Qwen3.6-27B com 0,751 e pelo Gemma-4-E4B com 0,618. A comparação por código revela diferenças e empates; as conclusões se restringem à amostra avaliada.
\end{abstract}

\section{Objetivo e escopo}
A análise compara a detecção dos problemas previamente catalogados. Não avalia a qualidade textual das explicações, recomendações, edições de código, acessibilidade da escrita ou efeito de RAG. Os resultados não representam uma avaliação geral dos modelos.

\section{Dados e método}
O conjunto contém dez códigos, dois por cenário, em cinco grupos: responsabilidades e coesão (C1), acoplamento concreto (C2), encapsulamento de regras e estado (C3), dificuldade de extensão (C4) e estado global ou dependências ocultas (C5). Há 24 problemas esperados no total. Cada código foi analisado por três modelos em cinco repetições, totalizando 150 execuções.

Cada problema possui um identificador, como \texttt{C1-A-01}. Os IDs listados em "Problemas detectados" representam problemas previstos identificados corretamente. Cada ID único conta como verdadeiro positivo (VP); a quantidade de IDs deve corresponder à coluna VP. O total esperado por código é VP+FN.

\[
\mathrm{Precisão}=\frac{VP}{VP+FP},\qquad
\mathrm{Revocação}=\frac{VP}{VP+FN},\qquad
F_1=\frac{2VP}{2VP+FP+FN}.
\]
As contagens são somadas por repetição antes do cálculo das métricas, conforme o plano de avaliação. A tabela geral apresenta a média das cinco repetições completas; para F1, também é apresentado o desvio padrão amostral.

\section{Resultados gerais}
\begin{table}[ht]
\centering
\caption{Contagens acumuladas em 50 execuções e médias das métricas entre cinco repetições.}
\begin{tabular}{lrrrrrr}
\toprule
Modelo & VP & FP & FN & Precisão & Revocação & $F_1$ médio $\pm$ DP\\
\midrule
M1 -- Gemma-4-E4B & 73 & 43 & 47 & 0,629 & 0,608 & 0,618 $\pm$ 0,031\\
M2 -- Qwen3.6-27B & 86 & 23 & 34 & 0,790 & 0,717 & 0,751 $\pm$ 0,044\\
M3 -- GPT-5.6 Luna & 104 & 18 & 16 & 0,854 & 0,867 & 0,860 $\pm$ 0,026\\
\bottomrule
\end{tabular}
\end{table}

M3 teve o maior F1 e revocação médios, mais acertos e menos falsos positivos. M2 ficou em posição intermediária. M1 apresentou menor F1 e mais falsos positivos. O desvio padrão observado de F1 foi menor para M3 (0,026), seguido por M1 (0,031) e M2 (0,044). Com cinco repetições, essa comparação é descritiva.

\section{Acertos por código}
A Tabela \ref{tab:acertos} mostra os IDs corretamente identificados nas cinco repetições. O denominador é cinco vezes a quantidade de problemas esperados no código. Como o número de repetições é igual para todos, o total acumulado e a média de acertos por execução produzem a mesma liderança. Empates foram mantidos.

\begin{table}[ht]
\centering
\caption{Acertos acumulados por código em cinco repetições}
\label{tab:acertos}
\scriptsize
\begin{tabular}{lrrrrl}
\toprule
Código & Esperados & M1 & M2 & M3 & Líder(es)\\
\midrule
C1-A & 3 & 6/15 & 6/15 & 11/15 & M3\\
C1-B & 2 & 9/10 & 9/10 & 10/10 & M3\\
C2-A & 2 & 4/10 & 6/10 & 4/10 & M2\\
C2-B & 2 & 10/10 & 10/10 & 10/10 & M1, M2, M3\\
C3-A & 3 & 5/15 & 15/15 & 15/15 & M2, M3\\
C3-B & 3 & 0/15 & 3/15 & 14/15 & M3\\
C4-A & 2 & 10/10 & 10/10 & 8/10 & M1, M2\\
C4-B & 2 & 10/10 & 7/10 & 7/10 & M1\\
C5-A & 3 & 10/15 & 10/15 & 15/15 & M3\\
C5-B & 2 & 9/10 & 10/10 & 10/10 & M2, M3\\
\bottomrule
\end{tabular}
\end{table}

\begin{figure}[ht]
\centering
\begin{tikzpicture}
\begin{axis}[
 ybar,bar width=5pt,width=\textwidth,height=7cm,ymin=0,ymax=16,
 ylabel={IDs corretos acumulados em cinco repetições},xlabel={Código},
 symbolic x coords={C1-A,C1-B,C2-A,C2-B,C3-A,C3-B,C4-A,C4-B,C5-A,C5-B},
 xtick=data,x tick label style={rotate=35,anchor=east,font=\small},
 ymajorgrids=true,grid style={gray!25},enlarge x limits=0.06,
 legend style={at={(0.5,1.03)},anchor=south,legend columns=3}]
\addplot+[fill=blue!55,draw=blue!70!black] coordinates {(C1-A,6)(C1-B,9)(C2-A,4)(C2-B,10)(C3-A,5)(C3-B,0)(C4-A,10)(C4-B,10)(C5-A,10)(C5-B,9)};
\addplot+[fill=orange!65,draw=orange!80!black] coordinates {(C1-A,6)(C1-B,9)(C2-A,6)(C2-B,10)(C3-A,15)(C3-B,3)(C4-A,10)(C4-B,7)(C5-A,10)(C5-B,10)};
\addplot+[fill=green!55,draw=green!70!black] coordinates {(C1-A,11)(C1-B,10)(C2-A,4)(C2-B,10)(C3-A,15)(C3-B,14)(C4-A,8)(C4-B,7)(C5-A,15)(C5-B,10)};
\legend{M1 -- Gemma,M2 -- Qwen,M3 -- GPT}
\end{axis}
\end{tikzpicture}
\caption{Acertos por código. O máximo possível é 10 ou 15 conforme a quantidade de problemas esperados.}
\end{figure}

M3 liderou sozinho em C1-A, C1-B, C3-B e C5-A e empatou em C3-A e C5-B. M2 liderou C2-A; M1 liderou C4-B. C3-B apresentou a maior diferença: M1 identificou 0/15 problemas possíveis nas cinco repetições, M2 3/15 e M3 14/15.

\section{Questões de pesquisa}
\begin{description}
\item[QP1 -- Eficácia geral.] F1 médio: M3 (0,860), M2 (0,751) e M1 (0,618), restrito a esta amostra.
\item[QP2 -- Falsos positivos.] Foram 18 para M3, 23 para M2 e 43 para M1 nas 50 execuções de cada modelo.
\item[QP3 -- Tipo de problema.] A liderança varia por código; a média geral não mostra sozinha essa variação e os empates.
\item[QP4 -- Consistência.] O DP amostral de F1 foi 0,026 para M3, 0,031 para M1 e 0,044 para M2. Cinco repetições não bastam para inferir estabilidade ampla.
\end{description}

\section{Limitações e conclusão}
A amostra possui dez códigos construídos, e todos contêm problemas previstos. A análise assume corretas as classificações humanas de VP, FP e FN. Não mede desempenho em códigos sem problemas, qualidade textual das explicações, recomendações, edições ou efeito de RAG. Não foram realizados testes inferenciais.

Nos dados disponíveis, M3 teve o maior desempenho geral. Por código, M1 ou M2 igualaram ou superaram M3 em alguns casos. Para encontrar o modelo com mais acertos, contam-se os IDs corretos de cada código ou suas somas nas cinco repetições, mantendo os empates.

\section*{Fontes}
\begin{enumerate}
\item \textit{ATLAS -- Plano de Avaliação}, inventário de códigos e problemas previstos.
\item \textit{Execucoes\_organizadas\_M1\_M2\_M3\_v2.xlsx}, aba "Execuções".
\end{enumerate}
\end{document}
